\documentclass[12pt]{article}

\usepackage{amsmath, amssymb, amsthm}
\usepackage{booktabs}
\usepackage{geometry}
\usepackage{hyperref}
\usepackage{graphicx}
\usepackage{xcolor}

\geometry{margin=1in}

\newtheorem{definition}{Definition}
\newtheorem{theorem}{Theorem}
\newtheorem{remark}{Remark}

\title{\textbf{Flood Forecasting as a Least Squares Problem}\\
\large 18.0851: Computational Science & Engineering}
\author{Rutvik Deshpande}
\date{}

\begin{document}

\maketitle
% ============================================================
\section{Problem Setup}

A flood wave entering a river takes days to travel downstream. If we can observe water levels at upstream gauges today, we have a mathematical signal about what will happen at a downstream location in the future. The question is: how do we turn those observations into a prediction?

We have $n$ days of historical data: daily streamflow at a target gauge and at several upstream gauges, plus precipitation measurements. The goal is to find a linear function of today's observations that best predicts flow at the target gauge $k$ days from now using least squares. Matrix $A$ is built from observations, a vector $b$ from future flow values, to solve;

\[
\min_{x \in \mathbb{R}^p} \|Ax - b\|_2^2.
\]
The solution $\hat{x}$ is a vector of weights. Applied to any new day's observations, it gives a forecast.

% ============================================================
\section{Building the Design Matrix}

\subsection{Features and Normalization}

On each day $t$, various set of measurements are observed. After preprocessing, these become a feature vector $a_t \in \mathbb{R}^p$. The features used are:
\begin{itemize}
    \item Log-transformed flow at the target gauge: $\log(1 + Q_{\text{target}}(t))$
    \item Daily and rolling precipitation sums (7-day, 14-day)
    \item Log-transformed flow at each upstream gauge: $\log(1 + Q_u(t))$ for $u \in \mathcal{U}$
    \item Sub-catchment precipitation between consecutive gauge pairs
\end{itemize}

The streamflow values were log-transformed because daily discharge is highly skewed and can vary across several orders of magnitude. This transformation reduces the influence of extreme flood peaks, preventing a small number of large events from dominating the squared-error loss and making the model training more stable.

All features are z-score normalized using training-set statistics. For each feature $j$:
\[
\tilde{a}_{tj} = \frac{a_{tj} - \mu_j}{\sigma_j}
\]
where $\mu_j$ and $\sigma_j$ are the training mean and standard deviation. This ensures every feature lives on the same scale before forming the matrix product $A^T A$.

\subsection{The Design Matrix}

Stack $n$ normalized feature vectors as rows, and prepend a column of ones for the intercept:
\[
A = \begin{bmatrix}
1 & \tilde{a}_{1,1} & \tilde{a}_{1,2} & \cdots & \tilde{a}_{1,p-1} \\
1 & \tilde{a}_{2,1} & \tilde{a}_{2,2} & \cdots & \tilde{a}_{2,p-1} \\
\vdots & \vdots & \vdots & \ddots & \vdots \\
1 & \tilde{a}_{n,1} & \tilde{a}_{n,2} & \cdots & \tilde{a}_{n,p-1}
\end{bmatrix}
\in \mathbb{R}^{n \times p}
\]

The target vector collects the log-flow $k$ days ahead:
\[
b = \begin{bmatrix} \log(1 + Q_{\text{target}}(t_1 + k)) \\ \log(1 + Q_{\text{target}}(t_2 + k)) \\ \vdots \\ \log(1 + Q_{\text{target}}(t_n + k)) \end{bmatrix} \in \mathbb{R}^n
\]

In this project, four separate least-squares systems are constructed, one for each forecast horizon 
\(k \in \{1,2,3,4\}\). The feature matrix \(A\) remains fixed across horizons, while the target vector 
\(b\) changes depending on how many days ahead the model is asked to predict.

For the Tar River dataset, the training set contains \(n = 7{,}178\) daily observations and 
\(p = 26\) model parameters, corresponding to 25 input features plus an intercept term. Since the 
number of observations is much larger than the number of unknown parameters, the system is 
overdetermined. In practice, this means there is no exact solution that satisfies every daily 
observation perfectly, so the model estimates the set of coefficients that best fits the data in 
the least-squares sense.
% ============================================================
\section{Least Squares}

We have far more equations than unknowns: the design matrix $A$ has $n \approx 7{,}000$ rows (one per training day) and only $p = 26$ columns (one per feature). No vector $x$ exists that satisfies $Ax = b$ exactly. Some residual is unavoidable. The question becomes which choice of $x$ leaves the smallest residual, and how we should measure "smallest."

We measure it by the sum of squared errors. For each training day $i$, the model produces a prediction $a_i^T x$, and we compare it to the observed value $b_i$. Squaring the per-day error and summing across days gives:
\[
f(x) = \|Ax - b\|_2^2 = (Ax - b)^T(Ax - b)
\]
Squaring is a deliberate choice. It penalizes large errors disproportionately, which matters here: missing a flood event by 5{,}000 cfs is much worse than ten ordinary days each off by 500 cfs. It also makes $f(x)$ smooth and differentiable, so we can find the minimum with calculus rather than search.

Expanding the product gives a clean quadratic in $x$:
\[
f(x) = x^T A^T A x - 2 b^T A x + b^T b
\]
The function is a paraboloid bowl: positive semi-definite quadratic, single minimum, no local traps. The minimum is found by setting the gradient to zero:
\[
\nabla_x f(x) = 2 A^T A x - 2 A^T b = 0
\]
which rearranges into the \textbf{normal equations}:
\[
\boxed{A^T A \, \hat{x} = A^T b}
\]

This is the central equation of the project. It converts an unsolvable system of 7{,}000 equations into a solvable system of 26. The solution $\hat{x}$ is the vector of feature weights that minimizes total squared prediction error across the entire training record.